\documentclass{article}
\usepackage{amsmath} % para ambientes matemáticos

\begin{document}

A matemática está presente em várias situações do dia a dia, desde calcular o troco até estimar distâncias e medidas.  
Ela nos ajuda a entender padrões e relações entre quantidades.

Podemos representar operações simples utilizando símbolos matemáticos. Por exemplo, a soma de dois números $a$ e $b$ pode ser escrita como $a + b$, e uma fração  como $\frac{a + c}{b}$.

Outra forma de mostrar operações é usando o ambiente de equação para destacá-las:

\begin{equation}
x_1 + y_1 = 10
\end{equation}

\begin{equation}
2 \div z = 8
\end{equation}

Além disso, podemos representar frações, potências e raízes de maneira clara:

\begin{equation}
\frac{a}{b},\quad x^2, \quad \sqrt{y}
\end{equation}

A matemática não serve apenas para cálculos abstratos; ela aparece em situações concretas, como:

\begin{enumerate}
    \item Calcular a  média das notas escolares.  
    \item Medir a quantidade de ingredientes em uma receita.  
    \item Estimar a distância entre dois pontos em um mapa.
\end{enumerate}

Outras aplicações importantes incluem:

\begin{itemize}
    \item Determinar uma sequência usando $u_{n+1}=2u_{n}+3$.  
    \item Calcular juros simples com $J = P \times i \times n$.\footnote{Na fórmula de juros simples, \(P\) representa o principal, \(i\) a taxa e \(n\) o número de períodos.}.
    \item Resolver problemas de geometria, como encontrar a área da hipotenusa $c = \sqrt{a^2 + b^2}$.  
\end{itemize}

\end{document}
